\chapter{Related work\label{cha:chapter3}}

In this chapter, we discuss the relevant literature review for the chosen topic.

Identify the important subcomponents (e.g., in terms of methods, evaluation, preprocessing or visualization) of your thesis and for each subcomponent, make a separate section with subsections as shown below.

The summary of the literature is shown in Table~\ref{}.



\section{Satellite SIF Enhancement Techniques}
\label{sec:sif_enhancement}

Satellite observations of solar-induced chlorophyll fluorescence (SIF) provide information on vegetation photosynthetic activity, but instrument characteristics restrict their direct use. Sensors with relatively small footprints, such as OCO-2 and OCO-3, sample narrow and discontinuous ground tracks. Wide-swath instruments such as TROPOMI and GOME-2 provide more complete coverage at coarser spatial support. SIF enhancement methods address these limitations using spatially complete environmental variables, relationships among SIF observations, or improved retrieval from measured radiance.

The literature can be divided by its primary objective. \emph{Reconstruction and gap filling} estimate SIF at unobserved locations or dates and can extend a short record using an earlier predictor archive. \emph{Spatial downscaling} estimates how coarse observed or reconstructed SIF is distributed among finer subpixels. Unconstrained methods apply a coarse statistical relationship to fine predictors, whereas aggregation-constrained methods preserve the parent mean or total. A third group contains retrieval enhancement, forecasting, and hybrid geostatistical methods. The groups overlap because a product may reconstruct and downscale simultaneously or use machine learning within an observation-preserving redistribution.

This distinction is important when interpreting spatial resolution. A value reported on a 500~m grid is not necessarily a SIF observation at 500~m: fine variation is inferred from assumptions and auxiliary variables. Reaggregation to the supervisory grid demonstrates coarse consistency but not individual fine-pixel accuracy. The review therefore considers model and output characteristics together with whether evaluation uses held-out supervisory data, independent SIF observations, or proxies such as gross primary production (GPP), crop yield, and drought indices.

\subsection{Reconstruction and Gap Filling}

Early globally contiguous products established the general data-driven reconstruction strategy. GOSIF aggregates quality-controlled OCO-2 SIF and learns its relationship with MODIS enhanced vegetation index (EVI), photosynthetically active radiation (PAR), vapour-pressure deficit, and air temperature using Cubist model trees. Applying the model to complete gridded predictors fills the OCO-2 swath gaps and extends the record to 2000 \cite{li2019global}. Its temporally held-out test result showed that sparse OCO-2 observations can support a global reconstruction, while the lower performance for some biomes indicated that one global score does not describe all vegetation regimes. CSIF uses a simpler feed-forward neural network with four MODIS reflectance bands to reconstruct clear-sky instantaneous and all-sky daily SIF \cite{zhang2018global}. Training on 2015--2016 and evaluation on 2014 and 2017 provided temporal separation between model fitting and testing. Because CSIF principally represents the optical canopy state and illumination, the residual between observed OCO-2 SIF and CSIF can retain short-term environmental down-regulation that is not visible in reflectance alone.

Several studies modify this universal mapping to account for ecological or observational context. The SIFoco2\_005 product fits separate artificial neural networks (ANNs) for each biome and 16-day time step, using seven-band MODIS reflectance to fill OCO-2 orbital gaps \cite{yu2019high}. Geographic sub-biomes improve the distribution of training samples, while time-specific models allow the SIF--reflectance relationship to vary with phenology and stress. Removing the biome and time stratification caused seasonal under- and overestimation and reduced drought sensitivity. Independent airborne measurements from the Chlorophyll Fluorescence Imaging Spectrometer (CFIS), including observations away from OCO-2 tracks, provided stronger evidence than a same-sensor random split. The spatially continuous TanSat product similarly uses a random forest and MODIS optical variables, but explicitly demonstrates the importance of the cosine of solar zenith angle as a proxy for illumination \cite{ma2020generation}. Its reduced accuracy after omitting this variable shows that observation geometry and radiation conditions should be represented when SIF samples span changing illumination.

Reconstruction can also extend the temporal range of a short, high-quality sensor record. RTSIF trains XGBoost on clear-sky TROPOMI SIF with MODIS reflectance, land-surface temperature and land cover, CERES PAR, and C3/C4 vegetation fractions, and then applies the relationship to the longer ancillary archive \cite{chen2022long}. The variables are motivated through a light-use-efficiency interpretation: reflectance characterizes canopy state and absorbed radiation, while temperature and vegetation type help represent changes in fluorescence efficiency. DOSIF pursues a still finer temporal objective by producing daily, globally contiguous OCO-3-like SIF \cite{yu2025dosif}. Its Moving Spatial--Temporal Window Sampling strategy constructs a separate model for each day of year from a temporal window across multiple years and geographic sub-biomes. CatBoost is selected after comparison with other algorithms. This contextual sampling outperforms a universal model and emphasizes that SIF--reflectance relationships vary with region and phenological stage. Independent CFIS evaluation strengthens the product assessment, although performance values reported in different parts of the paper are not fully consistent and should be interpreted using the detailed results.

ST-LGBM addresses the more demanding problem of extrapolating into complete regions between OCO-2 swaths \cite{shen2022spatiotemporal}. In addition to NIRv, PAR, temperature, vapour-pressure deficit, and land cover, it derives a spatial SIF factor from geographically nearby observations with similar vegetation state and a temporal factor from corresponding phenological periods in other years. These observation-derived constraints allow LightGBM to use information from neighbouring swaths rather than relying only on coincident ancillary variables. The study is particularly informative because it compares random validation with a swath-held-out experiment. The ordinary LightGBM test performance declined more strongly when complete regions were withheld, whereas ST-LGBM was more robust. This demonstrates why randomly scattered test samples can overestimate the ability to fill genuinely unobserved spatial gaps.

Regional reconstruction can target information that is lost in coarser temporal products. The ROSIF study fits separate XGBoost models between eight GOCI reflectance bands and OCO-3 SIF near 10:00 and 14:00 local time, producing continuous 500~m morning and afternoon fields \cite{he2025tracking}. SIF is subsequently normalized using PAR and NIRv to approximate fluorescence yield, and the morning-to-afternoon difference is used as a drought indicator. The approach shows that water and heat stress may appear in sub-daily photosynthetic depression rather than in a single observation. However, the physiological interpretation relies on simplifying assumptions in the NIRv-based yield decomposition, particularly for sparse vegetation. Together, these reconstruction studies show a progression from universal optical mappings toward context-specific, observation-constrained, and sub-daily models. They also show that model performance depends on the validation geometry and on whether the covariates can represent rapid physiological responses rather than only structural vegetation change.

\subsection{Spatial Downscaling}

Statistical downscaling commonly assumes that a relationship learned at coarse support can be evaluated using the same predictors at finer resolution. In Henan Province, a random forest is trained between monthly GOSIF, NDVI, and daytime land-surface temperature at the GOSIF grid and then applied to 1~km predictors \cite{zhang2024spatial}. The enhanced field supports a SIF anomaly drought index and shows stronger association with MODIS GPP than the original product. A related study uses a CNN with NDVI, EVI, daytime temperature, and nighttime temperature to downscale GOSIF to approximately 1~km \cite{zhang2020downscaling}. Convolutional filters incorporate adjacent-pixel context, which is potentially useful in fragmented farmland. In both studies, however, evaluation relies mainly on reaggregated agreement with GOSIF and relationships with GPP, drought, or yield. These results demonstrate application value, but they do not constitute independent measurements of SIF at the new fine scale. Moreover, because GOSIF is itself reconstructed, its upstream assumptions may propagate into the second downscaling stage.

Aggregation-constrained approaches retain more information from the original SIF observation. The global GOME-2 method of Duveiller et al. uses a semi-empirical light-use-efficiency formulation to distribute coarse SIF according to fine-resolution NIRv, NDWI, and afternoon land-surface temperature \cite{duveiller2020spatially}. These variables represent absorbed-light and canopy effects, water limitation, and temperature stress. OCO-2 is used to select the configuration, and a pixel-wise correction harmonizes the product with TROPOMI. Unlike direct statistical prediction, the framework preserves the parent GOME-2 signal while estimating its internal distribution. It therefore provides an early precedent for respecting aggregation support, although the subpixel pattern remains conditional on the selected proxies and semi-empirical assumptions.

Bias and residual correction provide an intermediate strategy between direct prediction and strict redistribution. BCSIF first uses a random forest to predict fine-resolution TROPOMI SIF from optical, illumination, and environmental variables, then estimates the discrepancy between observed and predicted SIF at the coarse scale and distributes this bias to the 0.005-degree grid \cite{hu2024spatially}. The correction is intended to restore physiological information present in TROPOMI but missing from the universal predictor relationship. Its validation at two ChinaSpec sites suggests that reducing spatial mismatch can improve agreement with tower SIF, although the limited number of sites constrains generalization. DSIFRFK\_EA0.05 likewise decomposes coarse GOME SIF into a random-forest trend and residuals that are interpolated by ordinary kriging before being added to the fine prediction \cite{jin2024downscaling}. The use of AVHRR and ERA5 enables reconstruction before the MODIS era. Nevertheless, moving residuals from 2-degree support to a $0.05^{\circ}$ grid does not eliminate the change-of-support problem; some apparent fine detail necessarily comes from interpolation assumptions.

Deep-learning methods use spatial neighbourhoods to represent heterogeneity within a coarse SIF cell. SIFnet learns a tenfold scale transformation by coarsening TROPOMI SIF to $0.5^{\circ}$ and training a residual CNN to recover the original $0.05^{\circ}$ field from coarse SIF and finer auxiliary rasters \cite{gensheimer2022convolutional}. The learned transformation is then applied to native TROPOMI and 0.005-degree covariates to produce approximately 500~m SIF. Its loss combines mean squared error with structural dissimilarity, thereby evaluating both magnitude and spatial pattern. Training outside North America and holding North America out for validation provides a geographic transfer test, while averaging the 500~m predictions within independent OCO-2 and OCO-3 footprints directly resembles a footprint-aware validation design. The parent TROPOMI SIF is the most important feature and NIRv the strongest auxiliary variable. However, the loss does not explicitly require the fine-cell mean to equal its parent SIF observation, and performance remains heterogeneous in drylands and complex urban landscapes.

TroDSIF makes the observation-preservation step explicit. A random forest first estimates 500~m SIF, but those predictions serve only as relative spatial weights in a Gaussian redistribution of original TROPOMI SIF \cite{chen2025improved}. Consequently, model-generated spatial detail is controlled by the contemporaneous coarse measurement rather than replacing it. Reaggregation closely reproduces the original TROPOMI signal, and tower measurements from six ChinaSpec sites provide independent fine-scale evidence. The principal limitation is the assumed linear proportionality between the original SIF and the predicted SIF used as a redistribution weight. If this relationship fails, the method can preserve the parent value while allocating it incorrectly among subpixels.

HSIF provides a physically constrained alternative. It requires the mean of 1~km subpixels to reproduce the low-resolution TROPOMI observation and uses a light-use-efficiency and radiative-transfer formulation to separate directional observed SIF from total canopy fluorescence emission \cite{zhang2023generating}. Chlorophyll FPAR, fluorescence efficiency, and escape probability govern the allocation. The method therefore avoids learning an unrestricted statistical scale transformation and explicitly conserves energy. Its comparison with flux-tower GPP further shows that agreement is strongest when the SIF averaging radius approaches the tower footprint. However, HSIF cannot reconstruct dates without an original TROPOMI observation and is sensitive to uncertainty in BRDF, FPAR, and escape-probability inputs.

CNSIF combines long-term reconstruction with CNN downscaling and is particularly relevant to fragmented agricultural landscapes \cite{du2025cnsif}. A 2D CNN receives GOSIF together with NIRv, EVI, near-infrared reflectance, and land-surface temperature. Neighbourhood windows represent the fine spatial structure within a coarse cell, and transfer learning applies a model developed at coarse support to 500~m Landsat/Sentinel-2 reflectance and Landsat thermal data. The resulting maps improve the delineation of winter-wheat fields, urban boundaries, and heterogeneous ecosystems. Yet the method inherits uncertainty from GOSIF, and its held-out 2019 comparison assesses consistency with the supervisory product rather than absolute fine-pixel accuracy. Independent tower SIF is therefore the more informative evaluation. Overall, the downscaling literature demonstrates that finer auxiliary imagery can add spatial structure, but also that coarse conservation, predictor transferability, and validation support determine whether that structure can be interpreted as enhanced SIF rather than merely a sharpened proxy field.

\subsection{Other and Hybrid Enhancement Approaches}

Some studies improve SIF in ways other than reconstructing or downscaling observations for the same date. Pais et al. first fill gaps in OCO-2 SIF with a random forest and then forecast monthly and seasonal SIF using two earlier SIF maps, EVI, and spatial and temporal variables \cite{pais2025forecasting}. They compare several statistical and deep-learning models, but errors from the first gap-filling step can affect the later forecast. Li et al. instead improve the original SIF retrieval by averaging matched OCO-2/3 observations within TROPOMI footprints and using them to train ANNs on TROPOMI radiance spectra \cite{li2025more}. This approach reduces noise in individual SIF retrievals and can provide better training data for later reconstruction or downscaling. Hybrid geostatistical methods use both predictor variables and information from nearby SIF observations. Tadi\'c et al. use CSIF to represent the overall spatial pattern and nearby OCO-2 observations to correct it locally, although their leave-one-out validation may be easier than filling a large gap between satellite swaths \cite{tadic2024hybrid}. coSIF jointly models monthly OCO-2 SIF and carbon dioxide for the following month, allowing cokriging to estimate both missing SIF and its uncertainty; however, when the relationship with carbon dioxide is weak, a simpler SIF-only kriging model performs similarly \cite{jacobson2023spatial}.

\subsection{Synthesis}

No single strategy resolves all limitations of satellite SIF. Universal reconstructions provide continuous records but may smooth biome-specific physiology. Unconstrained downscaling produces fine patterns, but coarse consistency and GPP correlation do not establish fine-pixel accuracy. Observation-preserving and physical approaches protect the parent signal, although subpixel allocation still depends on model assumptions. Geostatistical methods retain local observation structure and quantify uncertainty, but weaken as gaps grow and nearby observations disappear.

Three recurring research needs follow from this literature. First, spatial or temporal holdouts are necessary to test genuine extrapolation rather than interpolation among neighbouring samples. Second, external SIF measurements from another satellite, aircraft, or tower should be distinguished from proxy validation against GPP, yield, or drought indicators. Third, fine-resolution models should make the spatial support of their supervision and evaluation explicit. Few existing studies combine high-resolution Sentinel-2 predictors, the actual OCO-2 footprint support, agricultural composition information, comparisons between tabular and spatial neural-network inputs, and a direct analysis of how aggregating the supervisory target changes stability and predictive performance. These gaps motivate the enhancement experiments developed in this thesis.

\section{Remote Sensing-Based Crop Yield and GPP Estimation}
\label{sec:crop_yield_gpp}

Crop yield results from photosynthetic carbon uptake, respiration, biomass allocation, and the effects of weather, soil, management, and stress throughout the growing season. Remote sensing observes only parts of this process, but it provides repeated and spatially extensive information about canopy condition and productivity. The reviewed studies broadly use four approaches: SIF-based yield estimation, optical and spectral-index models, multi-source prediction, and SIF-based GPP estimation. These groups overlap because several yield models first estimate GPP, while many optical models also include weather or crop-model variables.

Data-driven SIF studies relate seasonal fluorescence directly to reported yield or include it in machine-learning models with vegetation, climate, and soil variables \cite{peng2020assessing}\cite{zheng2025modeling}\cite{joshi2023winter}. Peng et al. compared gap-filled OCO-2, TROPOMI, and coarse GOME-2 SIF with MODIS indices and land-surface temperature for maize and soybean in the U.S. Midwest. Higher-resolution OCO-2 and TROPOMI SIF outperformed GOME-2 and gave the best forward predictions for 2018. However, NIRv sometimes matched or exceeded SIF, and performance depended on crop type, validation method, and training-record length. For winter wheat in the North China Plain, SIF-based XGBoost achieved an independent 2020 result of $R^2=0.87$, while a CONUS study found that adding SIF to NIRv increased explained yield variation from 64\% to 69\%. Collectively, these studies show that SIF adds photosynthetic information, but partly overlaps with optical and environmental predictors and is not automatically the strongest variable in every setting.

SIF is particularly informative when stress changes photosynthesis before visible canopy structure. During the 2012 U.S. Midwest drought, yield declined by 25\% and GOSIF by 22\%, whereas NDVI, EVI, and NIRv declined by only 4--10\%; GOSIF also had the strongest yield relationship ($R^2=0.91$) \cite{qiu2022monitoring}. During the 2010 Indian heatwave, SIF detected reduced winter-wheat activity in early March, before NDVI and EVI responded \cite{song2018satellite}. Fluorescence yield decreased more than fPAR, indicating reduced photosynthetic efficiency rather than only less absorbed light. Nevertheless, SIF predicted a 13.9\% yield loss compared with the reported loss of approximately 6\%, so early stress detection does not guarantee equally accurate final-yield estimation.

Process-guided SIF models instead represent the steps from fluorescence to GPP and harvested yield \cite{guan2016improving}\cite{he2020ground}\cite{kira2024scalable}\cite{xue2025lightweight}\cite{wang2025practical}. Common components include corrections for canopy escape probability, differences between C$_3$ and C$_4$ photosynthesis, respiration, the NPP:GPP ratio, harvest index, and harvested area. He et al. found that seasonal TROPOMI SIF explained tower GPP with $R^2=0.89$, while accounting for planted area and C$_3$/C$_4$ composition increased the county productivity relationship from $R^2=0.72$ to 0.86. Kira et al.'s mechanistic light-reaction model was comparable to ANN and random forest for U.S. corn but outperformed them for wheat in the data-limited Indo-Gangetic Plain, showing the value of process knowledge when ground data are scarce. More recent lightweight models estimate the fraction of open PSII reaction centres, $q_L$, from SIF and environmental stress. Wang et al. reproduced leaf-level $q_L$ even under high light or temperature and obtained an average county-yield $R^2$ of 0.78 for U.S. corn and soybean. Such models are more interpretable than direct regressions, but still depend on crop-specific parameters and assumptions used to convert photosynthetic uptake into yield.

The choice between data-driven and process-guided models is therefore mainly a trade-off between flexibility and transferability. Machine-learning models can learn complex regional relationships when long, reliable yield records are available, but can also learn local correlations that fail in a new year or region. Process-guided models require assumptions about fluorescence, respiration, and crop allocation, yet need fewer yield observations and provide a clearer explanation of how environmental stress affects production. Hybrid approaches that retain these physical relationships while calibrating a limited number of parameters may be most useful where ground data are scarce.

Optical approaches estimate yield from canopy greenness, chlorophyll, water status, and structure using indices or high-resolution imagery \cite{vannoppen2021estimating}\cite{chandel2019yield}\cite{hunt2019high}\cite{vallentin2022suitability}. Their performance strongly depends on phenological timing and evaluation scale. In Belgium, seasonal NDVI summaries were weak winter-wheat predictors, while precipitation near tillering and anthesis produced $R^2=0.66$. Conversely, heading-stage NDVI and NDWI explained approximately 95--96\% of grain-yield variation in a controlled Indian nitrogen experiment, but without independent regional validation. Sentinel-2 enables operationally relevant spatial detail: Hunt et al. mapped U.K. wheat yield at 10~m and reduced RMSE from 0.66 to 0.61~t~ha$^{-1}$ by adding environmental data. However, the model used local combine-harvester data from one year. A 13-year northeast German analysis likewise found correlations ranging from almost zero to 0.94, with a median of only 0.19. Higher resolution and red-edge bands were useful, but correct acquisition timing and field heterogeneity were often more important than the choice of index.

These contrasting results show why an index should not be evaluated independently of the crop calendar. A measurement near heading or grain filling can contain more yield information than a seasonal maximum, while the same index may saturate in a dense canopy or respond too slowly to short stress events. High spatial resolution improves field delineation, but it cannot compensate for imagery acquired at an uninformative growth stage.

Multi-source studies combine satellite observations with weather, soil, phenology, or crop simulations because no single measurement describes the complete yield-formation process \cite{goihl2024crop}\cite{marszalek2022prediction}\cite{brandt2024ensemble}\cite{zhao2020predicting}\cite{srivastava2022winter}. Their main advantage is complementary information: optical data describe the canopy, weather and soil describe growth constraints, and crop models summarize processes such as water stress. In southern Germany, raw Sentinel-2 bands combined with precipitation and evapotranspiration explained 84\% of winter-wheat yield variance. In Australia, adding a crop-model water-stress index to Sentinel-2 structural and chlorophyll indices produced an independent-test $R^2=0.93$. Brandt et al. combined six estimators for approximately 140,000--155,000 German parcels per year; majority voting achieved winter-wheat $R^2=0.74$ at parcel level and 0.79--0.89 after district aggregation. This improvement partly reflects spatial averaging and may hide local errors. Validation design is therefore essential: Goihl found that regional random-forest models often transferred better than one Saxony-wide model, while strong training fits sometimes failed in the withheld year. Srivastava et al. similarly withheld complete years and found that a one-dimensional CNN learning weekly environmental sequences reduced RMSE by 7--14\% relative to the best baseline. These studies show the value of local context and temporal structure, but also the risk of overfitting heterogeneous regions.

Performance values from these studies are not directly interchangeable because their prediction targets range from field and parcel yields to county or district averages. Randomly divided samples can share the same year, region, and weather conditions, whereas a withheld year or region tests genuine extrapolation. Aggregation also reduces local noise and usually increases $R^2$. Consequently, the geographical and temporal separation of the test data is as important as the reported metric when judging whether a method could support an operational forecast.

The remaining studies focus on GPP, which is not harvested yield but provides the physiological link between SIF and biomass production \cite{hu2020detecting}\cite{ma2025gpp}\cite{shekhar2022well}. Downscaled GOME-2 SIF reproduced regional and seasonal GPP more effectively when local moving windows replaced one global regression, although its fine spatial pattern partly came from MODIS predictors. Transfer learning offers another solution to scale mismatch: Ma et al. pretrained a model on spatially extensive SIF-derived GPP and fine-tuned it with accurate but sparse eddy-covariance observations, improving spatial $R^2$ over both single-source alternatives. A comparison of four reconstructed SIF products found that they generally predicted tower GPP better than NDVI and EVI, with four-day CSIF and eight-day GOSIF performing best. However, the SIF--GPP relationship weakened at about 30\% of drought-affected sites, showing that product choice, temporal scale, land cover, and stress conditions influence reliability.

Accurate GPP estimation does not by itself determine grain yield. Respiration, the timing of carbon uptake, allocation to non-grain biomass, harvest index, and management all influence the final conversion. GPP validation is therefore strong evidence that a SIF product represents photosynthesis, but it remains proxy validation when the intended output is harvested crop yield.

Overall, SIF contributes direct information about photosynthetic activity and is especially useful for detecting rapid heat and drought responses, whereas Sentinel-2 and other optical sensors provide the spatial detail needed for field-level mapping. Weather, soil, phenology, crop composition, and process variables explain additional steps between canopy activity and harvested yield. The most convincing studies combine complementary data and use independent years, fields, regions, or flux measurements for evaluation. Common limitations remain coarse or mixed satellite footprints, limited ground observations, uncertain crop-area and yield statistics, and apparent accuracy gains caused by spatial aggregation. These findings support combining spatially enhanced SIF with high-resolution optical and environmental variables, while matching validation to the intended prediction scale.

\section{Comparative Analysis and Research Gap}
\label{sec:comparative_analysis_research_gap}

The selected enhancement methods differ mainly in their source of supervision, use of spatial context, and treatment of the original coarse SIF observation. GOSIF and CSIF reconstruct continuous global fields from sparse OCO-2 measurements using MODIS variables, but their pixel-wise models do not require the reconstructed values to preserve an observed coarse SIF total; GOSIF additionally includes radiation and meteorological variables, whereas CSIF relies mainly on reflectance \cite{li2019global,zhang2018global}. ST-LGBM also uses OCO-2 supervision, but adds spatial and temporal SIF factors to NIRv, PAR, temperature, vapour-pressure deficit, and land cover, making it better suited to large missing regions than a model based only on coincident predictors \cite{shen2022spatiotemporal}. BCSIF instead starts from spatially complete TROPOMI observations and uses a random forest with optical, illumination, and environmental variables, followed by a coarse-scale bias correction that reconnects the fine predictions to the measured TROPOMI signal \cite{hu2024spatially}. SIFnet uses TROPOMI and a residual CNN to learn spatial patterns from coarse SIF and fine auxiliary rasters, but it does not explicitly force the mean of its fine predictions to reproduce the parent observation \cite{gensheimer2022convolutional}. TroDSIF imposes this connection more directly by using random-forest predictions only as spatial weights for redistributing each contemporaneous TROPOMI observation, which improves coarse-scale consistency but still assumes that the predictor-based weights describe the correct subpixel pattern \cite{chen2025improved}. CNSIF uses a CNN and transfer learning with GOSIF, vegetation indices, near-infrared reflectance, and land-surface temperature, making it relevant to heterogeneous agricultural areas, although its 500~m output inherits uncertainty from the reconstructed GOSIF target \cite{du2025cnsif}. HSIF offers a physically constrained alternative that conserves the parent TROPOMI observation using absorbed light, fluorescence efficiency, and canopy escape probability, but it requires an observation on the target date and depends on uncertain physical inputs \cite{zhang2023generating}. Finally, coSIF uses monthly OCO-2 SIF, XCO$_2$, and spatial covariance in a cokriging framework that also quantifies uncertainty, but it becomes less useful when SIF--XCO$_2$ covariance is weak or nearby SIF observations are unavailable \cite{jacobson2023spatial}.

Across these approaches, the most important inputs are the parent SIF observation when one exists, optical variables such as reflectance and NIRv that describe canopy amount and absorbed light, and environmental variables such as PAR, temperature, vapour-pressure deficit, and soil moisture that describe illumination and physiological stress. Land cover and crop composition are also important in agricultural regions because a coarse satellite footprint can contain several crops and non-crop surfaces with different seasonal behaviour. Model architecture matters mainly through the form of information it can represent: tree models and feed-forward networks work well with tabular observations aggregated by pixel or footprint, whereas CNNs can additionally learn neighbourhood patterns, field boundaries, and spatial texture. However, a CNN cannot recover physiological information that is absent from its predictors, so predictor quality, spatial support, and validation design may be more important than choosing a convolutional rather than tabular model. A major research gap is that most widely evaluated products remain at approximately 500~m or coarser resolution, and independent SIF observations below 500~m are rarely available. Producing a 20~m SIF map and validating it after aggregation to 1~km can test whether the predictions reproduce SIF at the validation support, but it cannot by itself prove that every 20~m pixel is accurate. Gap filling OCO-2 is especially difficult because its narrow, separated swaths leave large areas without a contemporaneous parent SIF measurement, preventing the direct observation-preserving redistribution used by TROPOMI downscaling methods. This creates a need to test whether high-resolution Sentinel-2 bands and indices can provide useful sub-footprint information across missing OCO-2 regions, while recognizing that optical predictors may represent canopy structure more readily than rapid physiological changes. It also remains insufficiently tested whether explicit spatial models outperform tabular models when both use the same Sentinel-2 information and are evaluated with spatially separated samples. Since SIF has often represented photosynthetic stress and yield variation better than conventional greenness indices, isolating crop-pure SIF from mixed footprints may support more accurate crop-yield estimation \cite{qiu2022monitoring,song2018satellite,peng2020assessing}.
